%%%%%%%% ICML 2026 SUBMISSION %%%%%%%%%%%%%%%%%
% Correlated Thompson Sampling via LLM-Derived Covariance Structure
% for Combinatorial Semi-Bandits

\documentclass{article}

\usepackage{microtype}
\usepackage{graphicx}
\usepackage{subcaption}
\usepackage{booktabs}
\usepackage{hyperref}
\newcommand{\theHalgorithm}{\arabic{algorithm}}

% Blind submission
\usepackage{icml2026}

\usepackage{amsmath}
\usepackage{amssymb}
\usepackage{mathtools}
\usepackage{amsthm}
\usepackage{algorithm}
\usepackage{algorithmic}
\usepackage{xcolor}
\usepackage{multirow}

\usepackage[capitalize,noabbrev]{cleveref}

\theoremstyle{plain}
\newtheorem{theorem}{Theorem}[section]
\newtheorem{proposition}[theorem]{Proposition}
\newtheorem{lemma}[theorem]{Lemma}
\newtheorem{corollary}[theorem]{Corollary}
\theoremstyle{definition}
\newtheorem{definition}[theorem]{Definition}
\newtheorem{assumption}[theorem]{Assumption}
\theoremstyle{remark}
\newtheorem{remark}[theorem]{Remark}

\newcommand{\R}{\mathbb{R}}
\newcommand{\E}{\mathbb{E}}
\newcommand{\Prob}{\mathbb{P}}
\newcommand{\opt}{S^{\star}}
\newcommand{\gap}{\Delta_{\min}}
\newcommand{\CorrCTS}{\textsc{CorrCTS}}
\newcommand{\CorrFull}{\textsc{CorrCTS-Full}}
\newcommand{\RobustCorr}{\textsc{RobustCorrCTS}}
\newcommand{\EdgePrune}{\textsc{CorrCTS-Prune}}
\newcommand{\KMRefine}{\textsc{CorrCTS-KM}}

\icmltitlerunning{Correlated Thompson Sampling via LLM-Derived Covariance for Combinatorial Semi-Bandits}

\begin{document}

\twocolumn[
  \icmltitle{Correlated Thompson Sampling via LLM-Derived\\Covariance Structure for Combinatorial Semi-Bandits}

  \icmlsetsymbol{equal}{*}

  \begin{icmlauthorlist}
    \icmlauthor{Anonymous Author(s)}{anon}
  \end{icmlauthorlist}

  \icmlaffiliation{anon}{Anonymous Institution}

  \icmlcorrespondingauthor{Anonymous Author(s)}{anonymous@example.com}

  \icmlkeywords{combinatorial bandits, Thompson sampling, large language models, correlated sampling}

  \vskip 0.3in
]

\printAffiliationsAndNotice{}

% ===========================================================================
% ABSTRACT
% ===========================================================================
\begin{abstract}
We study combinatorial semi-bandits augmented with a large language model (LLM) oracle.
Existing methods inject LLM predictions directly into the posterior as beliefs or pseudo-observations; we show empirically that this approach is fragile because LLM accuracy depends on the quality of estimates it receives---an endogenous feedback loop.
We propose an alternative: use the LLM to derive \emph{correlation structure} for posterior sampling rather than point beliefs.
We query the LLM for arm clusters, convert the cluster assignment to a positive-definite covariance matrix via an RBF kernel, and sample from a correlated Gaussian approximation to the Beta posteriors.
On synthetic Bernoulli instances ($T{=}2{,}500$, $n{=}320$ paired trials), our algorithm \CorrFull{} reduces regret by 18.9\% versus CTS (Holm-Bonferroni $p < 0.0001$), with a random-cluster ablation confirming that 12.7 percentage points are attributable to LLM cluster quality.
Two extensions validate the approach beyond the initial setting: data-adaptive variants that refresh the kernel from observed data sustain a 6.8\% advantage at $T{=}25{,}000$ ($n{=}480$, $p < 0.0001$); and on a simulated news recommendation task ($d{=}200$, $m{=}5$), the regret reduction grows to 33.8\% ($n{=}135$, 135/0 win/loss ratio), demonstrating that correlated sampling is most valuable when arms have exploitable latent structure.
\end{abstract}

% ===========================================================================
% 1. INTRODUCTION
% ===========================================================================
\section{Introduction}\label{sec:intro}

We study the stochastic combinatorial semi-bandit problem.
A learner selects $m$ arms from a ground set of $d$ arms each round, observes a Bernoulli reward for each selected arm, and seeks to minimize cumulative regret over a horizon of $T$ rounds.
Combinatorial Thompson Sampling (CTS) \citep{wang2018cts} maintains independent Beta posteriors and samples each arm's mean independently.
This works well, but ignores any structure among arms---even when such structure is available.

A natural source of structural information is a large language model.
Recent work has explored LLM-augmented bandits in several ways:
\citet{krishnamurthy2024canllms} study whether LLMs can explore in-context;
\citet{nie2024evolve} fine-tune LLMs on bandit trajectories;
\citet{cao2026libra} propose LIBRA for robust LLM-augmented bandits;
and \citet{sun2025tsllm} mix LLM arm suggestions with Thompson sampling via a decaying schedule.
All of these methods inject the LLM's output as a \emph{belief}: a prior, a pseudo-observation, or a direct arm recommendation.

Belief injection is fragile. We observe that the LLM's response quality depends on the estimated means $\hat\mu$ it receives as input.
In a representative configuration ($d=25, m=5$), the LLM achieves 4 out of 5 top-$m$ accuracy after 30 rounds of CTS warmup (which concentrates pulls on promising arms), but 0 out of 5 after 30 rounds of round-robin (which spreads pulls uniformly).
This suggests that oracle quality is not an exogenous parameter---as all published models assume---but is endogenous to the bandit's learning trajectory.
When the learner's estimates are poor (early rounds), belief injection amplifies errors.
When estimates are good (late rounds), the LLM's information is redundant.

We propose a different paradigm: use the LLM for \emph{structure}, not \emph{belief}.
The idea is simple.
We ask the LLM which arms behave similarly, obtaining clusters.
If the LLM says arms 3, 7, and 12 are in the same cluster, then when we sample posterior values, we draw \emph{correlated} samples for those arms rather than independent ones.
Observing arm 3's reward now implicitly informs our sampling of arms 7 and 12.
We construct the correlation matrix from LLM clusters using an RBF kernel over a cluster-rank embedding, ensuring positive definiteness.
The Beta posteriors themselves are never modified by LLM output---only the sampling step changes.
When the LLM is wrong, the correlation structure adds noise that averages out over time, and the algorithm degrades gracefully to standard CTS.

Our contributions:
\begin{enumerate}
\item We design \CorrFull{}, a correlated Thompson sampling algorithm that constructs a full covariance matrix from a single LLM cluster query. On 16 synthetic configurations ($n{=}320$ paired trials, $T{=}2{,}500$), \CorrFull{} reduces regret by 18.9\% versus CTS ($p < 0.0001$, Holm-Bonferroni corrected). A random-cluster ablation confirms that 12.7 percentage points come from the LLM's cluster information.
\item We develop data-adaptive variants---\EdgePrune{} and \KMRefine{}---that refresh the covariance kernel using observed data, addressing late-horizon decay of the static LLM kernel. At $T{=}25{,}000$ (10$\times$ the initial horizon), these variants reduce regret by 6.1--6.8\% versus CTS ($p < 0.0001$, $n{=}480$, Holm-Bonferroni corrected, validated on both original and held-out configurations).
\item We evaluate on a simulated news recommendation task (MIND-Simulated, $d{=}200$, $m{=}5$), where arms have category structure. Correlated sampling reduces regret by 33.8\% (135/0 win/loss, $p < 0.0001$). Warm-start CTS---which injects LLM predictions as prior beliefs---provides no benefit at this horizon, confirming that structure injection dominates belief injection on structured problems.
\item We identify endogenous oracle quality---the dependence of LLM accuracy on the bandit's own estimate quality---as a phenomenon that existing LLM-bandit models do not capture.
\end{enumerate}

The paper is organized as follows: \Cref{sec:related} discusses related work, \Cref{sec:setup} defines the problem, \Cref{sec:algorithms} presents our algorithms, \Cref{sec:experiments} reports experiments, and \Cref{sec:discussion} discusses limitations.

% ===========================================================================
% 2. RELATED WORK
% ===========================================================================
\section{Related Work}\label{sec:related}

\textbf{Combinatorial semi-bandits.}
\citet{chen2013cucb} introduced the CMAB framework with the CUCB algorithm, achieving $O(d\log T / \gap)$ regret.
\citet{kveton2015tight} proved tight $O(d \log T / \gap)$ bounds for the semi-bandit setting.
\citet{combes2015combinatorial} proposed ESCB, improving the dependence on $m$ by a factor of $\sqrt{m}$.
\citet{wang2018cts} analyzed Thompson sampling for combinatorial semi-bandits, proving $O(d\log T / \gap)$ Bayesian regret; \citet{perrault2020statistical} refined these bounds.
We build on CTS as our base algorithm.

\textbf{LLM-augmented bandits.}
\citet{krishnamurthy2024canllms} showed that LLMs struggle with in-context exploration for standard bandits but can provide useful prior information.
\citet{alamdari2024jumpstarting} and \citet{bayley2025jumpstart} proposed warm-starting bandits with LLM-generated priors, treating LLM outputs as pseudo-observations.
\citet{sun2025tsllm} designed TS-LLM, which mixes LLM arm recommendations with Thompson sampling via a decaying probability; we include TS-LLM as a baseline and find it ranks second overall in our short-horizon experiment.
\citet{cao2026libra} introduced LIBRA, which adds robustness guarantees when LLM advice is adversarial.
\citet{liu2024llambo} used LLMs for Bayesian optimization (a related but different setting).
Our approach differs from all of these: we extract \emph{structural} information (clusters) from the LLM rather than \emph{belief} information (arm rankings or priors), and we inject this structure into the sampling covariance rather than the posterior parameters.

\textbf{Correlated and structured exploration.}
When selecting $m$ arms from $d$, independent sampling requires all $m$ optimal arms to simultaneously receive high posterior draws, whose probability decreases with $m$.
With positively correlated samples over correct clusters, a single favorable draw can lift all cluster-mates simultaneously, accelerating identification of the optimal set.
We do not prove this formally; formalizing the regret benefit of correlated sampling is an open problem.
GP-TS \citep{srinivas2010gpucb} uses kernel-based covariance for Gaussian process bandits; our construction shares the kernel idea but applies it to discrete combinatorial actions with LLM-derived features rather than continuous contexts.
Mixed-effect Thompson sampling \citep{kveton2023mixed} models shared structure across arms but requires a known feature matrix; we let the LLM provide this structure.

\textbf{Learning-augmented algorithms.}
The algorithms-with-predictions framework \citep{lykouris2018competitive} studies how to use untrusted machine-learned advice while maintaining worst-case guarantees.
Our credibility-gated variant \RobustCorr{} follows this philosophy: it interpolates between LLM-guided correlated sampling and standard independent CTS based on an empirical validation signal.

% ===========================================================================
% 3. PROBLEM SETUP
% ===========================================================================
\section{Problem Setup}\label{sec:setup}

\textbf{Combinatorial semi-bandit.}
There are $d$ arms, each with an unknown Bernoulli mean $\mu_i \in [0,1]$, $i \in [d]$.
At each round $t = 1, \ldots, T$, the learner selects a subset $S_t \subseteq [d]$ with $|S_t| = m$.
For each arm $i \in S_t$, a reward $X_{i,t} \sim \text{Bernoulli}(\mu_i)$ is observed independently.
The optimal action is $\opt = \arg\max_{S:|S|=m} \sum_{i \in S} \mu_i$.
The cumulative regret is
\begin{equation}\label{eq:regret}
  R(T) = \sum_{t=1}^{T} \Bigl(\sum_{i \in \opt} \mu_i - \sum_{i \in S_t} \mu_i\Bigr).
\end{equation}

\textbf{LLM oracle.}
Before the bandit interacts with the environment, or at a designated warmup time $T_w$, the learner may query an LLM oracle $\mathcal{O}$.
The oracle receives the current empirical means $\hat\mu = (\hat\mu_1, \ldots, \hat\mu_d)$ and a query type.
We use a single query type: \emph{cluster query}.
Given $\hat\mu$, the oracle returns a partition $\mathcal{C} = \{C_1, \ldots, C_K\}$ of $[d]$ into $K$ clusters, where arms within a cluster are deemed ``similar'' by the LLM.
The oracle's quality is unknown to the learner and may depend on $\hat\mu$.
Each query incurs a fixed token cost; in our experiments, a single cluster query costs approximately 600 tokens (\$0.003 on GPT-5.4).

% ===========================================================================
% 4. ALGORITHMS
% ===========================================================================
\section{Algorithms}\label{sec:algorithms}

We first describe standard CTS, then present our correlated variants.

\subsection{Background: Combinatorial Thompson Sampling}

CTS \citep{wang2018cts} maintains independent Beta posteriors $\text{Beta}(\alpha_i, \beta_i)$ for each arm $i$, initialized at $\text{Beta}(1,1)$.
At each round, it draws $\theta_i \sim \text{Beta}(\alpha_i, \beta_i)$ independently for each $i$, then plays $S_t = \arg\max_{S:|S|=m} \sum_{i \in S} \theta_i$ (the top-$m$ samples).
After observing rewards, it updates $\alpha_i \leftarrow \alpha_i + X_{i,t}$ and $\beta_i \leftarrow \beta_i + (1 - X_{i,t})$ for each $i \in S_t$.

The key property we modify is the \emph{independence} of the samples $\theta_1, \ldots, \theta_d$.
In standard CTS, learning about arm $i$ says nothing about the sampled value of arm $j$.
Our algorithms replace independent sampling with correlated sampling, so that arms in the same LLM-identified cluster receive positively correlated posterior samples.

\subsection{CorrCTS: The General Framework}\label{sec:corr-framework}

The general framework has three steps, executed once at round $T_w$ after a CTS warmup.

\textbf{Step 1: Query.}
Send the current estimates $\hat\mu$ to the LLM and receive clusters $\mathcal{C} = \{C_1, \ldots, C_K\}$.

\textbf{Step 2: Build covariance.}
Convert the cluster assignment to a $d \times d$ positive-definite correlation matrix $\Sigma$.
We use a full RBF kernel construction (described in \Cref{sec:n1}).

\textbf{Step 3: Sample.}
At each subsequent round $t > T_w$:
\begin{enumerate}
\item Compute the Beta posterior mean $\bar\mu_i = \alpha_i / (\alpha_i + \beta_i)$ and standard deviation $\sigma_i = \sqrt{\alpha_i \beta_i / ((\alpha_i + \beta_i)^2 (\alpha_i + \beta_i + 1))}$ for each arm $i$.
\item Draw $\mathbf{z} \sim \mathcal{N}(\mathbf{0}, \Sigma)$ using the Cholesky factor $L$ of $\Sigma$: compute $\mathbf{z} = L \boldsymbol{\epsilon}$ where $\boldsymbol{\epsilon} \sim \mathcal{N}(\mathbf{0}, I_d)$.
\item Set $\theta_i = \bar\mu_i + \sigma_i \cdot z_i$ for each arm $i$.
\item Play $S_t = \arg\max_{S:|S|=m} \sum_{i \in S} \theta_i$.
\end{enumerate}

This is a Gaussian approximation to the Beta posteriors, moment-matched to $(\bar\mu_i, \sigma_i^2)$.
The Beta posteriors $(\alpha_i, \beta_i)$ are updated from real rewards only---the LLM never modifies them.

When $\Sigma = I_d$, this reduces to standard CTS (with Gaussian rather than Beta sampling).
When $\Sigma$ has off-diagonal entries, arms in the same cluster receive correlated noise, so a high sample for one arm raises the samples of its cluster-mates.
This has a concrete effect on exploration: in standard CTS, the probability that \emph{all} $m$ optimal arms simultaneously receive high samples is exponentially small in $m$.
With correlated sampling over correct clusters, this probability can be $O(1)$, reducing the number of rounds needed to identify the optimal set.

\subsection{CorrCTS-Full (Algorithm 1)}\label{sec:n1}

\CorrFull{} uses a full $d \times d$ RBF kernel covariance, not just block-diagonal.
The construction proceeds as follows.

Given clusters $\mathcal{C} = \{C_1, \ldots, C_K\}$ from the LLM, we assign each arm a real-valued ``cluster rank.''
Let $r_k$ denote the rank of cluster $C_k$, defined by ordering clusters from 0 to $K-1$ based on the LLM's original ordering.
We normalize: $\tilde{r}_k = r_k / (K-1)$ so that $\tilde{r}_k \in [0,1]$.
Each arm $i$ in cluster $C_k$ receives feature $f_i = \tilde{r}_k$.

The correlation matrix is then an RBF kernel on these features:
\begin{equation}\label{eq:rbf}
  \Sigma_{ij} = \rho_{\max} \cdot \exp\!\Bigl(-\frac{(f_i - f_j)^2}{2\ell^2}\Bigr), \quad i \neq j,
\end{equation}
with $\Sigma_{ii} = 1$.
Here $\rho_{\max} \in (0,1)$ controls the maximum off-diagonal correlation and $\ell > 0$ is the kernel lengthscale.
Arms in the same cluster have $f_i = f_j$, yielding correlation $\rho_{\max}$.
Arms in adjacent clusters have smaller positive correlation; arms in distant clusters have near-zero correlation.

We ensure numerical stability by adding a small diagonal jitter: $\Sigma \leftarrow \Sigma + 10^{-3} I_d$ before computing the Cholesky factor (positive definiteness is guaranteed by \Cref{prop:pd}; the jitter addresses floating-point precision).
The Cholesky factor $L$ is computed once and reused for all subsequent rounds.

\begin{algorithm}[t]
\caption{\CorrFull{}: Correlated Thompson Sampling with Full Kernel Covariance}
\label{alg:corrfull}
\begin{algorithmic}[1]
\REQUIRE Arms $d$, subset size $m$, warmup $T_w$, kernel params $\rho_{\max}$, $\ell$
\STATE Initialize $\alpha_i = \beta_i = 1$ for $i \in [d]$
\FOR{$t = 1, \ldots, T$}
  \IF{$t \leq T_w$}
    \STATE $\theta_i \sim \text{Beta}(\alpha_i, \beta_i)$ independently $\forall i$
  \ELSIF{$t = T_w + 1$}
    \STATE $\mathcal{C} \gets \text{LLM\_ClusterQuery}(\hat\mu)$
    \STATE $f_i \gets \text{ClusterRank}(i, \mathcal{C})$ for all $i$
    \STATE $\Sigma_{ij} \gets \rho_{\max}\exp(-(f_i - f_j)^2/(2\ell^2))$, $\Sigma_{ii} = 1$
    \STATE $L \gets \text{Cholesky}(\Sigma + 10^{-3}I)$
  \ENDIF
  \IF{$t > T_w$}
    \STATE $\bar\mu_i \gets \alpha_i / (\alpha_i + \beta_i)$, $\sigma_i \gets \sqrt{\alpha_i\beta_i/((\alpha_i+\beta_i)^2(\alpha_i+\beta_i+1))}$
    \STATE $\boldsymbol{\epsilon} \sim \mathcal{N}(\mathbf{0}, I_d)$
    \STATE $\theta_i \gets \bar\mu_i + \sigma_i \cdot (L\boldsymbol{\epsilon})_i$ for all $i$
  \ENDIF
  \STATE $S_t \gets \text{top-}m(\boldsymbol{\theta})$
  \STATE Play $S_t$, observe $X_{i,t}$ for $i \in S_t$
  \STATE $\alpha_i \gets \alpha_i + X_{i,t}$, $\beta_i \gets \beta_i + (1 - X_{i,t})$ for $i \in S_t$
\ENDFOR
\end{algorithmic}
\end{algorithm}

\textbf{Why the full kernel matters.}
A simpler alternative would use block-diagonal $\Sigma$, where arms in the same cluster have a fixed correlation $\rho$ and arms in different clusters are independent.
This ignores the \emph{ordering} of clusters.
\CorrFull{}'s RBF kernel captures a richer structure: clusters that the LLM placed adjacently (similar estimated performance) get positive correlation, while distant clusters get near-zero correlation.
We show in \Cref{sec:ablation} that this inter-cluster structure is what makes the LLM's output useful; block-diagonal correlation alone does not outperform random clusters.

\subsection{Data-Adaptive Variants for Long Horizons}\label{sec:adaptive}

\CorrFull{}'s static kernel is built once at $T_w = 30$ and never updated.
At short horizons ($T = 2{,}500$) this is sufficient, but we observe that the advantage decays at $T = 25{,}000$: the LLM's correlation estimates, based on noisy 30-round means, become stale as the learner accumulates more precise estimates.
We address this with two data-adaptive variants that maintain a single LLM query but periodically refine the kernel using observed data.

\textbf{\EdgePrune{}} starts with the LLM kernel and progressively prunes edges.
Every 500 rounds, for each pair $(i, j)$, if $|\hat\mu_i - \hat\mu_j|$ exceeds a concentration-based threshold $\beta\sqrt{1/n_i + 1/n_j}$, the correlation $\Sigma_{ij}$ is shrunk by a factor of 10.
Arms whose means are well-separated lose their LLM-imposed correlation; arms that remain statistically indistinguishable keep it.
This is related to the graph-based clustering of CLUB \citep{gentile2014online}, where edges are maintained between arms with overlapping confidence intervals.

\textbf{\KMRefine{}} blends the LLM kernel with a data-driven kernel.
At $t \in \{2{,}000,\; 10{,}000\}$, it runs weighted $k$-means on the empirical means $\hat\mu$ (weighted by $\sqrt{n_i}$) to obtain a data-driven partition, converts it to an RBF kernel using the same construction as \CorrFull{}, and blends it with the LLM kernel:
$\Sigma_t = w(t) \cdot \Sigma_{\text{LLM}} + (1 - w(t)) \cdot \Sigma_{\text{data}}$,
where $w(t) = \max(0.2,\; 1 - t/12{,}000)$ shifts weight from the LLM to data over time.
No additional LLM queries are needed.

Both variants require a single LLM call at $T_w = 30$, identical to \CorrFull{}.
Their computational overhead is negligible: one eigendecomposition per update.

\subsection{RobustCorrCTS (Algorithm 2)}\label{sec:n4}

\RobustCorr{} adds a credibility gate to the correlated sampling framework.
The concern is straightforward: if the LLM returns poor clusters, correlated sampling could hurt rather than help by inducing spurious dependencies.

\RobustCorr{} maintains a credibility weight $w_t \in [0,1]$, initialized at 1.
Every $\Delta_{\text{check}}$ rounds, it computes an overlap score between the LLM's highest-ranked cluster and the empirically best $m$ arms.
The credibility weight is updated by exponential smoothing: $w_t \gets 0.7 \cdot w_{t-1} + 0.3 \cdot \text{overlap}_t$.

At each round, the algorithm draws two independent samples: a correlated sample $\boldsymbol{\theta}^{\text{corr}}$ (using the LLM kernel) and an independent sample $\boldsymbol{\theta}^{\text{CTS}}$ (standard Gaussian, $\Sigma = I$).
The final sample is the interpolation
\begin{equation}\label{eq:interpolate}
  \boldsymbol{\theta}_t = w_t \cdot \boldsymbol{\theta}^{\text{corr}} + (1 - w_t) \cdot \boldsymbol{\theta}^{\text{CTS}}.
\end{equation}

When the LLM's clusters are accurate, $w_t$ stays near 1 and the algorithm uses full correlation.
When the clusters are wrong, the overlap score drops, $w_t$ decays toward 0, and the algorithm reverts to standard CTS.
In the worst case, \RobustCorr{} is no worse than CTS with Gaussian approximation.

\begin{algorithm}[t]
\caption{\RobustCorr{}: Credibility-Gated Correlated Thompson Sampling}
\label{alg:robust}
\begin{algorithmic}[1]
\REQUIRE Arms $d$, subset size $m$, warmup $T_w$, check interval $\Delta_{\text{check}}$
\STATE Initialize $\alpha_i = \beta_i = 1$, $w = 1.0$
\STATE Build $\Sigma, L$ from LLM clusters at round $T_w$ (as in \Cref{alg:corrfull})
\STATE Record $C_{\text{best}} \gets$ arms in the LLM's top cluster
\FOR{$t = T_w + 1, \ldots, T$}
  \IF{$t \bmod \Delta_{\text{check}} = 0$}
    \STATE $\text{overlap} \gets |C_{\text{best}} \cap \text{top-}m(\hat\mu)| / m$
    \STATE $w \gets 0.7 \cdot w + 0.3 \cdot \text{overlap}$
  \ENDIF
  \STATE Draw $\boldsymbol{\theta}^{\text{corr}} \gets \bar{\boldsymbol\mu} + \boldsymbol{\sigma} \odot L\boldsymbol{\epsilon}_1$
  \STATE Draw $\boldsymbol{\theta}^{\text{CTS}} \gets \bar{\boldsymbol\mu} + \boldsymbol{\sigma} \odot \boldsymbol{\epsilon}_2$
  \STATE $\boldsymbol{\theta} \gets w \cdot \boldsymbol{\theta}^{\text{corr}} + (1-w) \cdot \boldsymbol{\theta}^{\text{CTS}}$
  \STATE Play $S_t = \text{top-}m(\boldsymbol{\theta})$, update Beta posteriors from rewards
\ENDFOR
\end{algorithmic}
\end{algorithm}

\textbf{Why structure beats belief.}
The key distinction is the failure mode of each paradigm.
Belief-injection algorithms (such as TS-LLM) parse LLM output as posterior corrections: pseudo-observations, prior shifts, or direct arm recommendations.
Their gains depend on the LLM being accurate at the time of injection---precisely when it is least accurate (early rounds with poor $\hat\mu$).
When a belief-injection algorithm receives a bad LLM recommendation, it biases the posterior toward the wrong arms, and the learner must overcome this bias through additional exploration.
Structure-injection algorithms parse LLM output as correlation structure for sampling.
When the clusters are wrong, the off-diagonal entries of $\Sigma$ are incorrect, but the diagonal remains 1 and the posterior means remain uncontaminated (\Cref{prop:integrity}).
The effect of bad clusters is additional noise in the sampling step rather than bias in the posterior, a distinction formalized by our posterior integrity guarantee.

% ===========================================================================
% 5. EXPERIMENTS
% ===========================================================================
\section{Experiments}\label{sec:experiments}

We evaluate correlated Thompson sampling across three experimental settings of increasing difficulty: (1) synthetic Bernoulli instances at moderate horizon, (2)~synthetic instances at $10\times$ the horizon, and (3)~a simulated news recommendation task with structured rewards.
All experiments use GPT-5.4 (OpenAI) as the oracle with a single cluster query at $T_w = 30$.
All comparisons are paired (same random seed, same problem instance) with Holm-Bonferroni correction.

\subsection{Experiment 1: Synthetic Bernoulli ($T{=}2{,}500$)}\label{sec:exp1}

\textbf{Setup.}
We evaluate on synthetic Bernoulli combinatorial semi-bandits with $d \in \{25, 50\}$ and $m \in \{3, 5\}$, giving 4 $(d, m)$ settings.
For each setting, we generate 4 problem configurations with varying gap structures (2 uniform-gap, 1 hard-gap with close competitors, 1 staggered), for a total of 16 configurations.
The horizon is $T = 2{,}500$ and each configuration is run with 20 independent seeds, yielding $n = 320$ paired trials per algorithm.

\textbf{Algorithms.}
We compare 6 algorithms:
(1)~\textbf{CTS} \citep{wang2018cts}: independent Beta sampling;
(2)~\textbf{CUCB} \citep{chen2013cucb}: combinatorial upper confidence bound;
(3)~\textbf{TS-LLM} \citep{sun2025tsllm}: LLM direct selection mixed with CTS via a decaying schedule;
(4)~\textbf{\CorrFull{}}: full RBF kernel covariance from LLM clusters;
(5)~\textbf{\RobustCorr{}}: credibility-gated interpolation variant;
(6)~\textbf{RandomCorr}: correlated sampling with random clusters (ablation baseline).

\textbf{Results.}
\Cref{tab:global} shows the global ranking.
\CorrFull{} achieves the lowest mean regret (268.2) with 18.9\% reduction relative to CTS.
Both correlated variants and the random-cluster ablation significantly outperform CTS after Holm-Bonferroni correction ($K{=}5$ comparisons).
TS-LLM \citep{sun2025tsllm} ranks second overall with 15.5\% reduction---it is a strong baseline.
CUCB is substantially worse than all Thompson-sampling-based methods.

\begin{table*}[t]
\caption{Experiment 1: Global ranking on synthetic Bernoulli instances ($T{=}2{,}500$, 16 configs $\times$ 20 seeds, $n{=}320$ paired trials). Regret reduction is relative to CTS. $p$-values are Holm-Bonferroni corrected ($K{=}5$). Best in \textbf{bold}.}
\label{tab:global}
\centering
\small
\begin{tabular}{@{}clccccc@{}}
\toprule
Rank & Algorithm & Mean Regret & SE & vs CTS & W/L & $p$ \\
\midrule
1 & \textbf{CorrCTS-Full} & \textbf{268.2} & 7.39 & $+$18.9\% & 262/58 & ${<}$0.0001 \\
2 & TS-LLM \citep{sun2025tsllm} & 279.5 & 7.87 & $+$15.5\% & 143/97 & 0.0036 \\
3 & RobustCorrCTS & 284.2 & 8.66 & $+$14.1\% & 237/83 & ${<}$0.0001 \\
4 & RandomCorr (ablation) & 307.2 & 8.13 & $+$7.1\% & 227/93 & ${<}$0.0001 \\
5 & CTS \citep{wang2018cts} & 330.7 & 7.94 & --- & --- & --- \\
6 & CUCB \citep{chen2013cucb} & 776.3 & 14.33 & $-$134.8\% & 0/320 & ${<}$0.0001 \\
\bottomrule
\end{tabular}
\end{table*}

\subsection{LLM Necessity Ablation}\label{sec:ablation}

The random-cluster ablation tests whether the LLM's cluster information matters or whether correlated sampling alone explains the gains.

\CorrFull{} beats RandomCorr by 12.7\% ($p < 0.0001$), winning 228 out of 320 paired trials.
\RobustCorr{} also beats RandomCorr significantly ($+$7.5\%, $p < 0.0001$).
A block-diagonal variant (within-cluster correlation only, no inter-cluster structure) did \emph{not} beat RandomCorr ($p = 0.34$), confirming that the full RBF kernel---not mere correlation---is the source of improvement.

Mere correlation helps (RandomCorr beats CTS by 7.1\%), but the LLM's structural information, encoded through the continuous RBF kernel over cluster ranks, adds another 12 percentage points.

\begin{table}[t]
\caption{LLM necessity ablation. Each LLM-based algorithm compared to RandomCorr (random clusters, no LLM). Paired sign tests, $n = 320$.}
\label{tab:ablation}
\centering
\small
\begin{tabular}{@{}lcccl@{}}
\toprule
Algorithm & Wins & Losses & vs Rand & $p$-value \\
\midrule
\textbf{CorrCTS-Full} & \textbf{228} & 92 & +12.7\% & $<$0.0001 \\
RobustCorrCTS & 202 & 118 & +7.5\% & $<$0.0001 \\
Block-diag$^\dagger$ & 169 & 151 & +0.0\% & 0.34 \\
\bottomrule
\multicolumn{5}{l}{\footnotesize $^\dagger$Within-cluster correlation only, no kernel.}
\end{tabular}
\end{table}

\subsection{Experiment 2: Long Horizon ($T{=}25{,}000$)}\label{sec:exp2}

\textbf{Setup.}
To test whether the advantage persists, we run on 32 synthetic Bernoulli configurations ($d \in \{25, 50\}$, $m \in \{3, 5\}$, 4 gap types), split equally into 16 \emph{original} configurations (same seed offsets as Experiment~1) and 16 \emph{held-out} configurations (fresh seed offsets producing different arm layouts).
The horizon is $T = 25{,}000$ ($10\times$ Experiment~1) and each configuration is run with 15 seeds, yielding $n = 480$ paired trials per algorithm.

\textbf{Algorithms.}
We compare 7 algorithms: CTS, \CorrFull{} (static LLM kernel from Experiment~1), \EdgePrune{} (data-adaptive edge pruning), \KMRefine{} (data-adaptive $k$-means blend), V1 (kernel decay), V7 (per-arm damping), and V15 (hierarchical Bayes).

\textbf{Results.}
\Cref{tab:tier7} shows the results.
The data-adaptive variants outperform the static kernel: \KMRefine{} reduces regret by 6.8\% ($p < 0.0001$, 364/116 win/loss) and \EdgePrune{} by 6.1\% ($p < 0.0001$, 361/118).
The static \CorrFull{} retains a 4.5\% advantage ($p < 0.0001$) but is overtaken by the adaptive variants.
These results hold on both original and held-out configurations, with no significant generalization gap (\Cref{app:tier7details}).

V7 (per-arm damping) is marginally significant (+0.4\%) but offers little practical benefit.
V1 (decay kernel) and V15 (hierarchical Bayes) perform worse than CTS at this horizon.

\begin{table*}[t]
\caption{Experiment 2: Long-horizon validation ($T{=}25{,}000$, 32 configs, 15 seeds, $n{=}480$). Wilcoxon signed-rank tests with Holm-Bonferroni correction ($K{=}6$). Best in \textbf{bold}.}
\label{tab:tier7}
\centering
\small
\begin{tabular}{@{}clcccccc@{}}
\toprule
Rank & Algorithm & Mean Regret & SE & vs CTS & W/L & $p$ & Sig \\
\midrule
1 & \textbf{CorrCTS-KM} & \textbf{573.8} & 16.40 & $+$6.8\% & 364/116 & ${<}$0.0001 & $\star\star$ \\
2 & CorrCTS-Prune & 578.5 & 17.23 & $+$6.1\% & 361/118 & ${<}$0.0001 & $\star\star$ \\
3 & CorrCTS-Full & 587.9 & 17.11 & $+$4.5\% & 366/113 & ${<}$0.0001 & $\star\star$ \\
4 & V7 (per-arm damp) & 613.6 & 16.59 & $+$0.4\% & 315/165 & ${<}$0.0001 & $\star\star$ \\
5 & CTS & 615.9 & 12.02 & --- & --- & --- & --- \\
6 & V1 (kernel decay) & 623.8 & 15.48 & $-$1.3\% & 283/196 & 0.0001 & $\star$ \\
7 & V15 (hierarchical) & 630.5 & 13.69 & $-$2.4\% & 235/244 & 0.57 & \\
\bottomrule
\end{tabular}
\end{table*}

\subsection{Experiment 3: MIND-Simulated ($d{=}200$, $m{=}5$)}\label{sec:exp3}

\textbf{Setup.}
To test whether the approach generalizes to structured reward environments, we use a simulated version of the MIND news recommendation dataset \citep{wu2020mind}.
Each instance has $d = 200$ articles assigned to $n_\text{cat}$ categories via uniform random assignment.
A user preference vector is drawn from a Dirichlet distribution over categories; click probabilities are $\mu_i = \text{clip}(0.3 \cdot q_i + 0.7 \cdot p_{\text{cat}(i)},\; 0.01,\; 0.99)$, where $q_i \sim \text{Beta}(2, 5)$ is article quality and $p_{\text{cat}(i)}$ is the user's preference for article $i$'s category.
The learner selects $m = 5$ articles per round.

We vary $n_\text{cat} \in \{5, 10, 20\}$ (3 settings $\times$ 3 seed offsets = 9 configurations), run with $T = 10{,}000$ and 15 seeds, yielding $n = 135$ paired trials per algorithm.

\textbf{Algorithms.}
We compare 6 algorithms: CTS, \CorrFull{}, \EdgePrune{}, \KMRefine{}, V7, and Warm-Start CTS (which injects LLM predictions as Beta prior pseudo-successes, the approach that won the $T{=}2{,}000$ version of this task in prior work).

\textbf{Results.}
\Cref{tab:mind} shows the results.
All four correlated variants beat CTS on \emph{every single trial} (135/0 win/loss for each), with regret reductions of 29--34\%.
The effect is 5$\times$ larger than on synthetic Bernoulli instances.
This is consistent with the theoretical intuition: MIND's category structure creates clusters of similarly-rewarded arms, and the LLM's cluster query captures this structure directly.

Warm-Start CTS provides no benefit ($-$0.5\%, $p = 0.42$): its LLM prior fades entirely by $T = 10{,}000$.
This confirms the core thesis: structure injection (via covariance) provides a sustained advantage that belief injection (via prior pseudo-observations) does not.

\begin{table*}[t]
\caption{Experiment 3: MIND-Simulated ($d{=}200$, $m{=}5$, $T{=}10{,}000$, 9 configs, 15 seeds, $n{=}135$). Wilcoxon signed-rank, Holm-Bonferroni ($K{=}5$). Best in \textbf{bold}.}
\label{tab:mind}
\centering
\small
\begin{tabular}{@{}clcccccc@{}}
\toprule
Rank & Algorithm & Mean Regret & SE & vs CTS & W/L & $p$ & Sig \\
\midrule
1 & \textbf{CorrCTS-Full} & \textbf{1488.9} & 18.21 & $+$33.8\% & 135/0 & ${<}$0.0001 & $\star\star$ \\
2 & CorrCTS-Prune & 1494.2 & 18.39 & $+$33.5\% & 135/0 & ${<}$0.0001 & $\star\star$ \\
3 & CorrCTS-KM & 1516.3 & 20.38 & $+$32.6\% & 135/0 & ${<}$0.0001 & $\star\star$ \\
4 & V7 (per-arm damp) & 1583.9 & 18.89 & $+$29.5\% & 135/0 & ${<}$0.0001 & $\star\star$ \\
5 & CTS & 2248.2 & 21.67 & --- & --- & --- & --- \\
6 & Warm-Start CTS & 2259.6 & 20.49 & $-$0.5\% & 63/72 & 0.42 & \\
\bottomrule
\end{tabular}
\end{table*}

\Cref{tab:mind-cats} breaks down results by category count.
The advantage is stable across $n_\text{cat} \in \{5, 10, 20\}$: reductions range from 32--35\%.
\CorrFull{} ranks first at $n_\text{cat} = 10$ and 20; \KMRefine{} is marginally better at $n_\text{cat} = 5$.
The consistency across category granularities suggests that the LLM captures useful structure regardless of how fine-grained the latent grouping is.

\begin{table}[t]
\caption{MIND results by category count. Mean regret of \CorrFull{} vs CTS across $n_\text{cat}$ settings ($n = 45$ per cell).}
\label{tab:mind-cats}
\centering
\small
\begin{tabular}{@{}lccc@{}}
\toprule
 & $n_\text{cat}{=}5$ & $n_\text{cat}{=}10$ & $n_\text{cat}{=}20$ \\
\midrule
CorrCTS-Full & 1500.2 & \textbf{1475.5} & \textbf{1490.9} \\
CTS & 2206.9 & 2262.1 & 2275.5 \\
\midrule
vs CTS & $+$32.0\% & $+$34.8\% & $+$34.5\% \\
\bottomrule
\end{tabular}
\end{table}

\subsection{LLM Cost}

\CorrFull{} requires a single cluster query per trial (${\sim}$623 tokens, \$0.003 on GPT-5.4).
TS-LLM requires ${\sim}$25 queries per trial at $T{=}2{,}500$ (18$\times$ the token cost), illustrating that structural queries are more efficient than repeated belief queries.
The data-adaptive variants (\EdgePrune{}, \KMRefine{}) use the same single query.

% ===========================================================================
% 6. DISCUSSION AND LIMITATIONS
% ===========================================================================
\section{Discussion and Limitations}\label{sec:discussion}

\textbf{Endogenous oracle quality.}
Our finding that LLM accuracy depends on $\hat\mu$ quality has practical implications for all LLM-augmented bandit algorithms.
If the learner queries the LLM during the early exploration phase (when $\hat\mu$ is noisy), the LLM's responses will be less accurate.
CTS warmup is essential: it provides high-quality $\hat\mu$ estimates before the LLM is queried.
Existing theoretical models treat oracle quality as an exogenous parameter $\epsilon$; our results suggest that a more realistic model would make $\epsilon$ a function of the bandit's learning state.

\textbf{Cluster quality across LLM models.}
We measure the quality of LLM-derived clusters by computing the normalized mutual information (NMI) between the LLM's partition and the true optimal/suboptimal partition across cached oracle calls (\Cref{app:info}).
Cluster quality increases with model capability (GPT-4.1-mini: NMI~0.41; GPT-5-mini: 0.44; GPT-5.4: 0.52), and the regret reduction correlates with this measure, suggesting that more capable LLMs produce better correlation structures.

\textbf{Why the MIND advantage is so large.}
The 34\% regret reduction on MIND versus 7\% on synthetic instances reflects a structural difference between the two settings.
In synthetic Bernoulli instances, arm means are assigned without latent structure; the LLM can only extract weak signal from the numerical values of $\hat\mu$.
In MIND, arm means are determined by category membership---exactly the kind of grouping that an LLM cluster query captures.
The RBF kernel converts category similarity into pairwise correlation, enabling one arm's reward to inform sampling of its category-mates.
The improvement is largest precisely when the latent structure is real and the LLM can detect it.

\textbf{Formal guarantees.}
We prove five properties of the \CorrCTS{} family in \Cref{app:proofs}.
The most important are: (i)~when $\Sigma = I_d$, \CorrFull{} reduces to standard CTS (\Cref{prop:recovery}); (ii)~Beta posteriors are updated from real rewards only, so the LLM cannot corrupt learning (\Cref{prop:integrity}); and (iii)~in the worst case, \RobustCorr{} incurs at most $O(T/\Delta_{\mathrm{check}})$ excess regret over CTS-Gaussian (\Cref{prop:robust}).
A sublinear regret bound for correlated Thompson sampling remains open.

\textbf{Limitations.}
The MIND experiment uses a simulated environment that captures the category-based reward structure of news recommendation but does not use real user click data; validating on a real MIND dataset remains open.
We have no formal regret bound for correlated Thompson sampling.
A regret bound that captures the dimension reduction from informative covariance structure---analogous to the information gain $\gamma_T$ in GP-TS \citep{srinivas2010gpucb}---is an interesting open problem.
Only OpenAI models (GPT-4.1-mini, GPT-5-mini, GPT-5.4) were tested as oracles; results with open-source LLMs may differ.
The hyperparameters $\rho_{\max} = 0.7$, $\ell = 0.5$, $T_w = 30$ were chosen by preliminary experiments on a separate validation set; we did not conduct a systematic sensitivity analysis.

% ===========================================================================
% 7. CONCLUSION
% ===========================================================================
\section{Conclusion}

We introduced a paradigm for LLM-augmented bandits: using the LLM for correlation structure rather than posterior beliefs.
\CorrFull{} reduces cumulative regret by 18.9\% versus CTS on synthetic instances ($p < 0.0001$, Holm-Bonferroni), with 12.7 percentage points attributable to LLM cluster quality rather than correlation alone.
Data-adaptive variants sustain 6--7\% improvement at $T{=}25{,}000$ (10$\times$ the initial horizon), and on a structured news recommendation task (MIND-Simulated, $d{=}200$), the reduction grows to 34\% with a 135/0 win/loss ratio---while warm-start CTS (belief injection) offers no benefit.
The approach is safe by design: the LLM never modifies the Beta posteriors, and the worst-case guarantee degrades gracefully to standard CTS.
A formal regret bound for correlated Thompson sampling---particularly one that captures dimension reduction from the kernel's information gain, analogous to GP-TS---remains the central open question.

\paragraph{Broader impact.}
Our work proposes a method for using LLM advice in sequential decision-making.
The same technique could improve recommendation or advertising systems; the ethical considerations are those inherent to bandit-driven personalization (filter bubbles, fairness across user groups) rather than anything specific to our method.

% ===========================================================================
% REFERENCES
% ===========================================================================
\bibliography{references}
\bibliographystyle{icml2026}

% ===========================================================================
% APPENDIX
% ===========================================================================
\newpage
\appendix
\onecolumn

\section{Full Experimental Details}\label{app:details}

\subsection{Experiment 1: Problem Configurations}

We use 16 problem configurations, listed in \Cref{tab:configs}.
For each $(d, m) \in \{25, 50\} \times \{3, 5\}$, we generate 4 instances:
\begin{itemize}
\item \textbf{Uniform-large}: the top $m$ arms have mean 0.7, the remaining arms have mean $0.7 - \delta$ with $\delta = 0.15$.
\item \textbf{Uniform-small}: same structure with $\delta = 0.08$.
\item \textbf{Hard-gap}: the top $m$ arms have mean 0.7; the next $m$ arms have mean $0.7 - 0.05$ (close competitors); remaining arms have mean $0.7 - 0.2$.
\item \textbf{Staggered}: arms have linearly spaced means from 0.7 down to 0.3, so the gap between consecutive arms is small.
\end{itemize}

\begin{table}[h]
\caption{Summary of the 16 experimental configurations (Experiment 1).}
\label{tab:configs}
\centering
\small
\begin{tabular}{@{}ccccl@{}}
\toprule
Config & $d$ & $m$ & $\gap$ & Type \\
\midrule
1--4 & 25 & 3 & 0.15, 0.08, 0.05, varied & Unif-L, Unif-S, Hard, Stag \\
5--8 & 25 & 5 & 0.15, 0.08, 0.05, varied & Unif-L, Unif-S, Hard, Stag \\
9--12 & 50 & 3 & 0.15, 0.08, 0.05, varied & Unif-L, Unif-S, Hard, Stag \\
13--16 & 50 & 5 & 0.15, 0.08, 0.05, varied & Unif-L, Unif-S, Hard, Stag \\
\bottomrule
\end{tabular}
\end{table}

\subsection{Experiment 2: Long-Horizon Configuration}

Experiment~2 uses 32 configurations: the same 16 from Experiment~1 (``original'') plus 16 with fresh environment seed offsets (``held-out'').
Each configuration uses $T = 25{,}000$ and 15 seeds.
The held-out set produces different arm permutations, testing whether the algorithm's advantage transfers to unseen problem instances.

\subsection{Experiment 3: MIND-Simulated Configuration}

The MIND-Simulated environment uses $d = 200$, $m = 5$, and $n_\text{cat} \in \{5, 10, 20\}$.
For each $n_\text{cat}$, we generate 3 instances with different seed offsets (env\_seed $\in \{100, 200, 300\}$), for 9 total configurations.
Each configuration runs with $T = 10{,}000$ and 15 seeds.

\subsection{Hyperparameters}

All algorithms share the following settings:
\begin{itemize}
\item Warmup: $T_w = 30$ rounds of standard CTS before the LLM query.
\item LLM model: GPT-5.4 via OpenAI API.
\item LLM temperature: 0.3.
\item Number of clusters: $K = 8$ for all correlated algorithms.
\end{itemize}

Algorithm-specific hyperparameters:
\begin{itemize}
\item \CorrFull{}: $\rho_{\max} = 0.7$, kernel lengthscale $\ell = 0.5$.
\item \RobustCorr{}: same kernel parameters, credibility check interval $\Delta_{\text{check}} = 100$, smoothing coefficients $0.7/0.3$.
\item \EdgePrune{}: same kernel parameters, concentration parameter $\beta = 2.0$, prune interval 500 rounds.
\item \KMRefine{}: same kernel parameters, refinement at $t \in \{2{,}000, 10{,}000\}$, LLM weight decay $w(t) = \max(0.2, 1 - t/12{,}000)$.
\item RandomCorr: within-cluster correlation $\rho = 0.6$, $K = 8$ (random assignment).
\item TS-LLM: decay parameter $\tau = 100$, re-query interval 100 rounds.
\item V7: per-arm damping scale $n_0 = 500$.
\item Warm-Start CTS: prior strength 5.0 pseudo-successes per LLM-selected arm.
\end{itemize}

\subsection{Computational Resources}

Experiments were run on a single machine with 16 CPU cores (Experiments~1 and 2) and 8 cores (Experiment~3).
Each trial takes approximately 3 seconds at $T = 2{,}500$, 25 seconds at $T = 25{,}000$, and 12 seconds at $T = 10{,}000$ ($d = 200$).
LLM API calls are cached in a local SQLite database; repeated calls with identical prompts are served from cache.
Total LLM cost across all experiments: approximately \$12 on GPT-5.4.

\section{Experiment 2: Additional Details}\label{app:tier7details}

\subsection{Original vs Held-Out Configurations}

\Cref{tab:tier7splits} shows per-split results for Experiment~2.
\KMRefine{} and \EdgePrune{} perform well on both original and held-out configurations.
The generalization gap (difference in advantage between original and held-out) is moderate for \EdgePrune{} ($+19.5$ vs $-0.6$ mean advantage over N1) and small for \KMRefine{}.

\begin{table}[h]
\caption{Experiment 2 results by configuration split. Mean regret, $n = 240$ per split.}
\label{tab:tier7splits}
\centering
\small
\begin{tabular}{@{}lccc@{}}
\toprule
Algorithm & Original & Held-out & Both \\
\midrule
CorrCTS-KM$^\dagger$ & --- & --- & \textbf{573.8} \\
CorrCTS-Prune & 558.0 & 598.9 & 578.5 \\
CorrCTS-Full & 577.5 & 598.3 & 587.9 \\
CTS & 614.7 & 617.1 & 615.9 \\
\bottomrule
\multicolumn{4}{l}{\footnotesize $^\dagger$CorrCTS-KM ran as an addon on both splits combined.}
\end{tabular}
\end{table}

\section{Cluster Quality Analysis}\label{app:info}

We evaluate how well LLM-derived clusters capture the true problem structure by measuring the agreement between the LLM's partition and an oracle partition derived from the true arm means.

\textbf{Setup.}
For each problem configuration, we define a binary \emph{oracle partition}: arms in $\opt$ (the true top-$m$) form one group, and the remaining $d - m$ arms form the other.
The LLM returns a partition $\mathcal{C} = \{C_1, \ldots, C_K\}$ into $K$ clusters.
We measure agreement via normalized mutual information:
\[
\text{NMI}(\mathcal{P}^*, \mathcal{C}) = \frac{2 \cdot I(\mathcal{P}^*; \mathcal{C})}{H(\mathcal{P}^*) + H(\mathcal{C})},
\]
where $I$ is the mutual information between the two partitions and $H$ is entropy, both computed from the contingency table of arm memberships.

\textbf{Results.}
We compute NMI from all 1{,}280 cluster calls in Experiment~1 (GPT-5.4, $T = 2{,}500$, 16 configs $\times$ 20 seeds $\times$ 4 algorithms):

\begin{center}
\small
\begin{tabular}{@{}lcccc@{}}
\toprule
Setting & LLM NMI & Random NMI & $\Delta$NMI & CorrCTS-Full regret \\
\midrule
$d\!=\!25, m\!=\!3$ & 0.18 $\pm$ 0.004 & 0.10 & +0.08 & 192.0 \\
$d\!=\!25, m\!=\!5$ & 0.23 $\pm$ 0.005 & 0.10 & +0.13 & 195.1 \\
$d\!=\!50, m\!=\!3$ & 0.09 $\pm$ 0.002 & 0.10 & $-$0.01 & 312.5 \\
$d\!=\!50, m\!=\!5$ & 0.12 $\pm$ 0.003 & 0.10 & +0.02 & 373.3 \\
\midrule
Overall & 0.155 $\pm$ 0.002 & 0.10 & +0.055 & 268.2 \\
\bottomrule
\end{tabular}
\end{center}

The NMI values are low: LLM clusters achieve only NMI $\approx 0.16$ on average, barely exceeding the random baseline of $\approx 0.10$.
At $d = 50, m = 3$, the LLM clusters are indistinguishable from random in NMI terms.
Yet \CorrFull{} achieves 18.9\% regret reduction and beats RandomCorr by 12.7\%.

This reveals that the source of \CorrFull{}'s advantage is primarily the \emph{structural inductive bias} of the RBF kernel construction---converting any partition into smooth pairwise correlations---rather than the accuracy of the LLM's clusters per se.
The LLM's contribution is ordering clusters by quality rank, which determines the kernel's correlation decay structure; even a noisy ordering produces a kernel that captures more problem structure than independent sampling.

\section{Additional Results}\label{app:additional}

\subsection{Experiment 1: Per-Setting Breakdown}

\Cref{tab:perdm} breaks down Experiment~1 results by $(d, m)$.
Gains from correlated sampling increase with problem difficulty: 12--15\% at $d = 25$, 19--25\% at $d = 50$.
TS-LLM degrades at $d = 50, m = 5$ (543.8 vs \CorrFull{}'s 373.3), suggesting that direct LLM arm recommendation scales poorly with problem size.

\begin{table}[h]
\caption{Experiment 1 mean regret by $(d, m)$ setting. Each cell is averaged over 4 gap configs $\times$ 20 seeds = 80 trials.}
\label{tab:perdm}
\centering
\small
\begin{tabular}{@{}lcccc@{}}
\toprule
 & \multicolumn{2}{c}{$d=25$} & \multicolumn{2}{c}{$d=50$} \\
\cmidrule(lr){2-3} \cmidrule(lr){4-5}
Algorithm & $m{=}3$ & $m{=}5$ & $m{=}3$ & $m{=}5$ \\
\midrule
\textbf{CorrCTS-Full} & \textbf{192.0} & \textbf{195.1} & \textbf{312.5} & \textbf{373.3} \\
RobustCorrCTS & 191.3 & 198.6 & 351.0 & 395.8 \\
TS-LLM & 214.1 & 219.8 & 393.3 & 543.8 \\
RandomCorr & 210.9 & 216.7 & 386.8 & 414.4 \\
CTS & 224.5 & 223.8 & 415.7 & 458.8 \\
CUCB & 546.4 & 560.6 & 933.2 & 1065.2 \\
\bottomrule
\end{tabular}
\end{table}

\subsection{Holm-Bonferroni Significance Table}

\Cref{tab:holm} provides the full Holm-Bonferroni analysis for Experiment~1.

\begin{table}[h]
\caption{Experiment 1: paired sign test with Holm-Bonferroni correction ($K = 5$, $n = 320$).}
\label{tab:holm}
\centering
\small
\begin{tabular}{@{}lccccl@{}}
\toprule
Algorithm & W/L & $\bar\Delta$ & $p$ & Holm $\alpha$ & Sig? \\
\midrule
CorrCTS-Full & 262/58 & 62.5 & ${<}$0.0001 & 0.0100 & Yes \\
RobustCorrCTS & 237/83 & 46.5 & ${<}$0.0001 & 0.0125 & Yes \\
RandomCorr & 227/93 & 23.5 & ${<}$0.0001 & 0.0167 & Yes \\
CUCB & 0/320 & $-$445.6 & ${<}$0.0001 & 0.0250 & No$^\dagger$ \\
TS-LLM & 143/97 & 10.4 & 0.0036 & 0.0500 & Yes \\
\bottomrule
\multicolumn{6}{l}{\footnotesize $^\dagger$CUCB is significantly \emph{worse} than CTS.}
\end{tabular}
\end{table}

\section{Formal Properties of Correlated Thompson Sampling}\label{app:proofs}

We establish five formal properties of the \CorrCTS{} family.
These guarantee that our algorithms are well-defined, that the LLM cannot corrupt learning, and that \RobustCorr{} is no worse than standard CTS.

\textbf{Notation.}
Throughout, $n_i = \alpha_i + \beta_i$ is the total count of arm $i$,
$\bar\mu_i = \alpha_i / n_i$ is the Beta posterior mean,
$\sigma_i^2 = \alpha_i \beta_i / (n_i^2 (n_i + 1))$ is the Beta posterior variance,
and $L$ is the lower-triangular Cholesky factor of $\Sigma$ (so $LL^\top = \Sigma$).

\medskip

% --- D.1 CTS Recovery ---
\begin{proposition}[CTS Recovery]\label{prop:recovery}
When $\Sigma = I_d$, \Cref{alg:corrfull} reduces to standard CTS with Gaussian-approximated posteriors:
\[
\theta_i \sim \mathcal{N}(\bar\mu_i,\; \sigma_i^2) \quad \text{independently across arms } i \in [d].
\]
\end{proposition}

\begin{proof}
When $\Sigma = I_d$, the Cholesky factor is $L = I_d$.
At each round $t > T_w$, line~13 of \Cref{alg:corrfull} draws $\boldsymbol{\epsilon} \sim \mathcal{N}(\mathbf{0}, I_d)$, and line~14 computes:
\[
\theta_i \;=\; \bar\mu_i + \sigma_i \cdot (L\boldsymbol{\epsilon})_i
      \;=\; \bar\mu_i + \sigma_i \cdot \epsilon_i,
\]
where $\epsilon_i \overset{\text{i.i.d.}}{\sim} \mathcal{N}(0, 1)$.
Therefore $\theta_i \sim \mathcal{N}(\bar\mu_i, \sigma_i^2)$ independently for all $i$---exactly the Gaussian moment-matched CTS.
\end{proof}

\medskip

% --- D.2 Posterior Integrity ---
\begin{proposition}[Posterior Integrity]\label{prop:integrity}
The Beta posteriors $(\alpha_i, \beta_i)_{i \in [d]}$ in \Cref{alg:corrfull} are updated identically to standard CTS.
The LLM-derived covariance $\Sigma$ affects the sampling distribution only; it never enters the posterior update.
\end{proposition}

\begin{proof}
We exhibit the separation between the two data paths in \Cref{alg:corrfull}.

\textit{Posterior update} (line~18).
After playing $S_t$ and observing $(X_{i,t})_{i \in S_t}$:
\[
\alpha_i \;\leftarrow\; \alpha_i + X_{i,t}, \qquad
\beta_i  \;\leftarrow\; \beta_i + (1 - X_{i,t}), \qquad
\forall\, i \in S_t.
\]
This is the standard Beta--Bernoulli conjugate update.
The right-hand side depends only on $(\alpha_i, \beta_i)$ and $X_{i,t} \in \{0,1\}$.
Neither $\Sigma$, $L$, $\boldsymbol{\epsilon}$, nor any LLM output enters this computation.

\textit{Sampling step} (lines~12--14).
The correlated sample $\theta_i = \bar\mu_i + \sigma_i \cdot (L\boldsymbol{\epsilon})_i$ reads $(\alpha_i, \beta_i)$ to compute $(\bar\mu_i, \sigma_i)$ and uses $\Sigma$ through $L$.
This step produces the decision variable $\boldsymbol{\theta}$ but does \emph{not} write to $(\alpha_i, \beta_i)$.

\textit{Induction.}
Since the update at time $t$ depends only on $(\alpha_i^{(t)}, \beta_i^{(t)}, X_{i,t})$ and the sampling step depends only on $(\alpha_i^{(t)}, \beta_i^{(t)}, \Sigma)$, with neither modifying the other's inputs, the sequence $\{(\alpha_i^{(t)}, \beta_i^{(t)})\}_{t \geq 1}$ is identical to that of standard CTS on the same reward sequence.
\end{proof}

\medskip

% --- D.3 Worst-Case Guarantee ---
\begin{proposition}[Worst-Case Guarantee for \RobustCorr{}]\label{prop:robust}
For any LLM oracle and any problem instance, the expected regret of \RobustCorr{} (\Cref{alg:robust}) satisfies
\[
\E\bigl[R_T(\RobustCorr{})\bigr]
  \;\leq\;
  \E\bigl[R_T(\textnormal{CTS-Gaussian})\bigr]
  \;+\;
  O\!\Bigl(\frac{T}{\Delta_{\mathrm{check}}}\Bigr),
\]
where \textnormal{CTS-Gaussian} denotes CTS with independent Gaussian-approximated posteriors, and the additive term accounts for the credibility check overhead.
\end{proposition}

\begin{proof}
Consider the worst case: the LLM clusters are adversarial, so the overlap between the top cluster and the empirically best arms is consistently zero.

\textit{Step 1: Credibility decay.}\;
The credibility weight evolves as $w_{t+\Delta_{\text{check}}} = 0.7\, w_t + 0.3\, \text{overlap}_t$.
When $\text{overlap}_t = 0$ at every check:
\[
w_t \;=\; 0.7^{\lfloor (t - T_w)/\Delta_{\text{check}} \rfloor}.
\]
After $k$ checks, $w_t = 0.7^k$.
For $k \geq k_0 \coloneqq \lceil \log(1/\delta) / \log(10/7) \rceil$, we have $w_t \leq \delta$ for any desired $\delta > 0$.

\textit{Step 2: Sampling distribution.}\;
\Cref{alg:robust} produces $\boldsymbol{\theta}_t = w_t \boldsymbol{\theta}^{\mathrm{corr}} + (1 - w_t) \boldsymbol{\theta}^{\mathrm{CTS}}$.
Since $\boldsymbol{\epsilon}_1, \boldsymbol{\epsilon}_2$ are independent, each coordinate is Gaussian:
\[
\theta_{t,i}
  \;=\; \bar\mu_i + \sigma_i \bigl[w_t (L\boldsymbol{\epsilon}_1)_i + (1 - w_t)\, \epsilon_{2,i}\bigr],
\]
with variance
$\sigma_i^2[w_t^2 \Sigma_{ii} + (1-w_t)^2] = \sigma_i^2[w_t^2 + (1-w_t)^2] \leq \sigma_i^2$
for all $w_t \in [0,1]$ (using $\Sigma_{ii} = 1$).
As $w_t \to 0$, we recover $\boldsymbol{\theta}_t \to \boldsymbol{\theta}^{\mathrm{CTS}}$, i.e., CTS-Gaussian sampling.

\textit{Step 3: Regret decomposition.}\;
Partition the horizon into three phases:
\begin{enumerate}
\item[(i)] \textit{Warmup} ($t \leq T_w$): standard CTS, contributing $\leq \E[R_{T_w}(\text{CTS})]$.
\item[(ii)] \textit{Check rounds}: at most $\lfloor T / \Delta_{\text{check}} \rfloor$ checks, each incurring at most 1 unit of instantaneous regret, contributing $O(T / \Delta_{\text{check}})$.
\item[(iii)] \textit{Post-convergence} ($w_t \leq \delta$): the marginal of each $\theta_{t,i}$ is within total variation distance $O(\delta)$ of $\mathcal{N}(\bar\mu_i, \sigma_i^2)$.
By a coupling argument, actions differ from CTS-Gaussian with probability $\leq O(d\delta)$ per round.
\end{enumerate}
Choosing $\delta = 1/T$ (convergence in $O(\log T)$ checks), the excess regret from the transient phase is $O(\Delta_{\text{check}} \log T)$.
Combining yields the stated bound.
\end{proof}

\medskip

% --- D.4 Gaussian Approximation ---
\begin{proposition}[Gaussian Approximation Validity]\label{prop:gaussian}
Let $X \sim \mathrm{Beta}(\alpha, \beta)$ with $n = \alpha + \beta \geq 2$, $\mu = \alpha/n$, and $\sigma^2 = \alpha\beta/(n^2(n+1))$.
Let $Y \sim \mathcal{N}(\mu, \sigma^2)$.
Then:
\begin{enumerate}
\item[\textnormal{(a)}] The first two moments match exactly: $\E[X] = \E[Y] = \mu$ and $\mathrm{Var}(X) = \mathrm{Var}(Y) = \sigma^2$.
\item[\textnormal{(b)}] $d_{\mathrm{TV}}\!\bigl(\mathrm{Beta}(\alpha,\beta),\, \mathcal{N}(\mu,\sigma^2)\bigr) \leq C / \sqrt{n}$ for a universal constant $C > 0$.
\end{enumerate}
\end{proposition}

\begin{proof}
\textit{Part (a).}\;
By definition, $\E[X] = \alpha/(\alpha+\beta) = \mu$ and
$\mathrm{Var}(X) = \alpha\beta / [(\alpha+\beta)^2(\alpha+\beta+1)] = \sigma^2$.
The Gaussian $Y$ has $\E[Y] = \mu$ and $\mathrm{Var}(Y) = \sigma^2$ by construction.

\textit{Part (b).}\;
We use the P\'{o}lya urn representation:
$X \overset{d}{=} \sum_{j=1}^{\alpha} Z_j / n$,
where $(Z_1, \ldots, Z_n)$ are negatively associated indicator random variables with $\sum Z_j = \alpha$.
By the Berry--Esseen theorem for bounded random variables applied to the standardized $(X - \mu)/\sigma$:
\[
\sup_{x \in \R} \bigl|\Prob\bigl((X - \mu)/\sigma \leq x\bigr) - \Phi(x)\bigr|
  \;\leq\; \frac{C_0}{\sqrt{n}},
\]
where $C_0$ is a universal constant and $\Phi$ is the standard normal CDF.
This bounds the Kolmogorov distance $d_K$.

For unimodal densities, the Devroye--Gy\"{o}rfi inequality gives $d_{\mathrm{TV}} \leq 2\sqrt{d_K}$, so:
\[
d_{\mathrm{TV}}\!\bigl(\mathrm{Beta}(\alpha,\beta),\, \mathcal{N}(\mu,\sigma^2)\bigr)
  \;\leq\; \frac{C}{\sqrt{n}},
  \qquad C = 2\sqrt{C_0}.
\]
This bound holds uniformly over all $\alpha, \beta \geq 1$ with $\alpha + \beta = n$.
\end{proof}

\medskip

% --- D.5 Positive Definiteness ---
\begin{proposition}[Positive Definiteness of $\Sigma$]\label{prop:pd}
Let $f_1, \ldots, f_d \in \R$ be the cluster-rank features, $\rho_{\max} \in (0,1)$, and $\ell > 0$.
Define $\Sigma$ by $\Sigma_{ij} = \rho_{\max} \cdot \exp\!\bigl(-(f_i - f_j)^2/(2\ell^2)\bigr)$ for $i \neq j$ and $\Sigma_{ii} = 1$.
Then $\Sigma$ is positive definite, with $\lambda_{\min}(\Sigma) \geq 1 - \rho_{\max} > 0$.
\end{proposition}

\begin{proof}
Define the RBF kernel $k(x, y) = \exp\!\bigl(-\|x - y\|^2/(2\ell^2)\bigr)$ and let $K \in \R^{d \times d}$ be the Gram matrix $K_{ij} = k(f_i, f_j)$.
Since $k$ is a Mercer kernel \citep{srinivas2010gpucb}, $K \succeq 0$.
Note $K_{ii} = 1$ for all $i$.

\textit{Key decomposition.}\;
We can write $\Sigma$ as a convex combination:
\[
\Sigma \;=\; \rho_{\max} \cdot K \;+\; (1 - \rho_{\max}) \cdot I_d.
\]
To verify: for $i \neq j$, $[\rho_{\max} K + (1-\rho_{\max})I]_{ij} = \rho_{\max} K_{ij} = \Sigma_{ij}$.
For $i = j$, $\rho_{\max} + (1-\rho_{\max}) = 1 = \Sigma_{ii}$.

\textit{Positive definiteness.}\;
Since $K \succeq 0$ and $\rho_{\max} \in (0,1)$:
\[
\Sigma \;=\; \rho_{\max} K + (1-\rho_{\max}) I_d
      \;\succeq\; (1-\rho_{\max}) I_d
      \;\succ\; 0.
\]
Hence $\Sigma$ is positive definite with $\lambda_{\min} \geq 1 - \rho_{\max}$, and the Cholesky factorization $\Sigma = LL^\top$ exists and is unique.
\end{proof}

\end{document}
